% Generated draft: requires mathematical and source review.
\documentclass[11pt]{article}
\usepackage[margin=25mm]{geometry}
\usepackage{fontspec}
\setmainfont{Latin Modern Roman}
\usepackage{amsmath,amssymb,mathtools,unicode-math}
\setmathfont{Latin Modern Math}
\usepackage{enumitem,needspace,xurl,hyperref}
\setlength{\parindent}{0pt}
\setlength{\parskip}{5pt}
\setlength{\emergencystretch}{3em}
\newcommand{\N}{\mathbb{N}}
\newcommand{\Z}{\mathbb{Z}}
\newcommand{\Q}{\mathbb{Q}}
\newcommand{\R}{\mathbb{R}}
\newcommand{\C}{\mathbb{C}}
\newcommand{\F}{\mathbb{F}}
\newcommand{\A}{\mathbb{A}}
\newcommand{\PP}{\mathbb{P}}
\newcommand{\E}{\mathbb{E}}
\newcommand{\eps}{\varepsilon}
\newcommand{\dd}{\,\mathrm{d}}
\begin{document}
\section*{UnsolvedMath ID: 20001095}
\textbf{Source problem code:} catalogue entry
\subsection{Definitions and mathematical statement}
R. Graham's source entry proposes studying weak \(k\)-term arithmetic progressions as a more flexible substitute for ordinary progressions, suggesting that standard problems might become easier. It does not state a particular extremal, density, Ramsey, or algorithmic conjecture. The stored delimiters appear as ``b'' and ``c''; the supplied aid reads them as floor brackets. Under that explicit interpretation a weak progression is
\[
 \{\lfloor i\alpha+\beta\rfloor:1\le i\le k\},
 \qquad \alpha\in\R,\ \alpha\ge1,\quad\beta\in\R.
\]
The floor convention, rather than ceiling or nearest-integer rounding, matters for the strict endpoint conditions below. The parameter bound ensures distinct increasing terms.

The aid's classification problem takes an increasing integer sequence \(a_1<\cdots<a_k\). It reports equivalence between being the displayed floor block and nonemptiness of
\[
 [1,\infty)\cap\bigcap_{1\le i<j\le k}
 \left(\frac{a_j-a_i-1}{j-i},\frac{a_j-a_i+1}{j-i}\right).
\]
For any \(\alpha\) in this intersection, the parameter recovery is
\[
 \max_i(a_i-i\alpha)\le\beta<\min_i(a_i+1-i\alpha).
\]
The aid reports that these \(O(k^2)\) exact rational inequalities construct admissible parameters. For \(k=1\), the empty intersection places no additional restriction on \(\alpha\ge1\), and the beta interval is nonempty.

Equivalently, its reported gap characterization is: there exists an integer \(q\ge1\) such that all gaps \(a_{i+1}-a_i\) lie in \(\{q,q+1\}\), and the binary word
\[
 w_i=a_{i+1}-a_i-q\quad(1\le i<k)
\]
is balanced. Here balanced means that any two contiguous subwords of equal length have numbers of ones differing by at most one. For \(k=1\), the gap condition is vacuous. The presence of two adjacent gap values alone is sufficient through \(k=4\), but not generally.

The aid gives the sharp first failure at \(k=5\):
\[
 \{0,1,2,4,6\},\qquad\text{gaps }1,1,2,2,\qquad w=0011.
\]
Its subwords \(00\) and \(11\) violate balance, and the strict slope intervals require both \(\alpha<3/2\) and \(\alpha>3/2\). The aid calls it a smallest counterexample to the gap-only criterion, meaning smallest number of terms, not a specified minimal diameter or lexicographic order. These are finite-block structural results; they do not settle any unspecified version of the standard additive-combinatorial problems mentioned programmatically by the source.
\subsection{Short English statement}
The source proposes replacing arithmetic progressions by finite floor-sequence blocks. The supplied aid characterizes such blocks using strict slope intervals and balanced adjacent-gap words, with a five-term example showing why two possible gap sizes alone are insufficient.
\subsection{Sources}
\begin{itemize}
\item Local UnsolvedMath registry, permanent ID \(20001095\), source code \texttt{catalogue entry}; \texttt{aim-combinatorics-notes.json}, \texttt{notes[219]}, tag \texttt{problem}.
\item AIM workshop ``Recent trends in additive combinatorics,'' Problem 1.14, R. Graham: \url{https://aimath.org/WWN/additivecomb/additivecomb.pdf}. Statement status \texttt{exact}, including the stored floor-delimiter artifacts. No PDF was accessed for this draft.
\item No original remarks, literature field, proposer/year metadata, or authoritative external citation for the finite floor-block characterization is supplied beyond the in-text R. Graham attribution.
\item Research attempt 1: \texttt{context\_only}, problem status at run \texttt{not\_a\_problem}, rigor label \texttt{proved}, novelty confidence low; registry status \texttt{open}, classification \texttt{OPEN-TRIAGE}, suggested difficulty L2. Created and updated 14 August 2026. This machine-generated aid is not peer reviewed; independently verify before citation.
\item Ulam AI / UnsolvedMath Contributors; CC BY 4.0.
\end{itemize}
\end{document}
